% ===================================================================
%  TABELLA N-RO 10 — Le mar. — Le porto.
%  Traduction 2026 d'apres texto/io, controlee sur texto/fr et
%  texto/es. Consigne : voir texto/ia/10-tabella-01.tex.
%
%  Le francais decale sa numerotation d'un cran a partir de l'alinea
%  8 ; c'est une coquille de l'atelier francais, et la colonne suit
%  l'ido, qui compte juste.
% ===================================================================

%%K t10-apar-1 apar
\VUsaut{4.0mm}
\VUtitre{15.0pt}{60}{TABELLA N\textsuperscript{o} 10}
\VUsaut{3.0mm}
\VUpk{11.4pt}{40}{Le mar. --- Le porto.}
\VUsaut{3.0mm}

%%K t10-01-1 p
1. --- In estate, durante que unes cerca calma e frescor sur un montania, alteres
prefere ir al litore. Assi nos faceva le anno passate, proque nos ama multo le
\VUgras{Oceano} \textsuperscript{(1)}.

%%K t10-02-1 p
2. --- Quanto a me, io senti un grande placer contemplante iste immense extension
de aqua que se extende usque le \VUgras{horizonte} \textsuperscript{(2)}, e
vidente partir pro un longe \VUgras{traversata} le enorme \VUgras{naves a
vapor} \textsuperscript{(3)} sur le quales imbarca le viageros qui va in
America.

%%K t10-03-1 p
3. --- Pro isto, mi patre, qui cognosce mi inclination, sempre elige, pro passar
le vacantias, un plagia multo calme, situate presso un de nostre grande
\VUgras{portos militar}.

%%K t10-03-2 p
Durante nostre ultime sojorno, le \VUgras{esquadra} veniva ancorar detra le
enorme \VUgras{dica} \textsuperscript{(4)} que claude e protege le \VUgras{porto
militar} \textsuperscript{(5)}, e io succedeva secundo mi desiro admirar le
diverse \VUgras{naves de guerra}, le parve \VUgras{submarino}
\textsuperscript{(10)} que plongia e dispare sub le aquas e anque le enorme
\VUgras{cuirassato} \textsuperscript{(9)} que, usque ora timeva quasi solmente
le attaccos del \VUgras{torpedinero} \textsuperscript{(8)}.

%%K t10-tit-2 sub
\VUsaut{3.0mm}
\VUpk{10.2pt}{40}{Le parve naves.}
\VUsaut{3.0mm}

%%K t10-04-1 p
4. --- Iste ultime ja sapeva multo \VUgras{habilemente} trovar le puncto debile
del spisse cuirasse de metallo pro fixar illac le periculose
\VUgras{torpedine}, del qual le explosion causa le perdita del nave, ma tamen
illo esseva minus timibile que le submarino, proque le contratorpedineros le
persequeva facilemente.

%%K t10-05-1 p
5. --- Io anque poteva vider le rapide \VUgras{crucciatores}
\textsuperscript{(6)} cuirassate, quasi tanto invulnerabile como le ver
cuirassato, plure \VUgras{naves de transporto} \textsuperscript{(12)} tres o
quatro \VUgras{cannoneras} \textsuperscript{(7)} e finalmente un
\VUgras{fregata} \textsuperscript{(11)}. \VUgras{Le spectaculo de iste flotta,
quando illo manovra pro levar le ancoras es le plus interessante.}

%%K t10-06-1 p
6. --- Quante differentia de construction inter iste grande massas de aciero e
le elegante \VUgras{tres-mastes} o le gratiose \VUgras{naves a vela}
\textsuperscript{(13)} ancora ora usate in le flotta commercial que io videva
entrar le porto, a tote velas, quando io promenava sur le \VUgras{quai}
\textsuperscript{(14)}. In ver istes perdeva multo de lor avantages
\VUgras{quando} le \VUgras{vento} esseva contrari. Illos debeva calar tote lor
velas pro reentrar le bassino e un \VUgras{remorcator} \textsuperscript{(16)}
esseva necessari pro cercar los e conducer los, como simple \VUgras{chalanas}
\textsuperscript{(17)}, al quai.

%%K t10-tit-3 sub
\VUsaut{3.0mm}
\VUpk{10.2pt}{40}{Le porto commercial.}
\VUsaut{3.0mm}

%%K t10-07-1 p
7. --- Lo que es interessante in le porto del qual io parla, es que illo contine
non solmente un porto militar, artificialmente arrangiate per gigantesc labores,
ma anque un meraviliose \VUgras{porto commercial} situate in le
\VUgras{estuario} de un grande fluvio.

%%K t10-08-1 p
8. --- Le lingua de terra que separa le duo bassinos es fortificate e defendite
per un serie de \VUgras{batterias} e \VUgras{bastiones} qui avantia sur le mar.
On ha besonio de un permission special pro promenar sur le \VUgras{remparos}
\textsuperscript{(15)}. Io lo obteneva facilemente, per mi oncle, qui es ingeniero
del \VUgras{cantieres} \textsuperscript{(18)} e io iva visitar le naves que on
construe sub vaste hangares.

%%K t10-09-1 p
9. --- Mesmo io assisteva al lanceamento de un cuirassato e io memora le
emotionante momento que tote le turba sentiva quando le grandissime massa,
abandonate a se mesme, comenciava glissar sur su quadro similante un cuna e
veniva flottar sur le aqua del fluvio.

%%K t10-10-1 p
10. --- Pro ir al cantieres, on transversa le \VUgras{campo de manovras}
\textsuperscript{(19)} ubi le soldatos del guarnition cata die veni exercer se,
on passa sur un \VUgras{ponticello} \textsuperscript{(20)} de ferro, que
communica le esplanada de parada con le \VUgras{arsenal}
\textsuperscript{(21)}.

%%K t10-tit-4 sub
\VUsaut{3.0mm}
\VUpk{10.2pt}{40}{Le arsenal e le dockos.}
\VUsaut{3.0mm}

%%K t10-11-1 p
11. --- Con mi oncle, io visitava iste grande construction plen de
\VUgras{cannones} \textsuperscript{(22)}, \VUgras{ballas}
\textsuperscript{(23)} e \VUgras{obuses}. Anque io videva le \VUgras{fabrica de
metallo} \textsuperscript{(24)} in le qual es fabricate le \VUgras{caldieras}
\textsuperscript{(25)} del naves a vapor.

%%K t10-12-1 p
12. --- Multo presso, il ha le \VUgras{dockos} \textsuperscript{(26)} --- dockos
de reparation --- e le multo remarcabile \VUgras{dockos flottante}
\textsuperscript{(27)} in le quales io poteva vider multo de presso, sur un nave
a reparar, le \VUgras{helice} \textsuperscript{(28)}, le \VUgras{quilia}
\textsuperscript{(29)} e le \VUgras{timon} \textsuperscript{(30)} de un barca.

%%K t10-13-1 p
13. --- Le porto commercial es tote tanto digne de studio como le porto militar,
e io passava multe horas admirante sur le quais ubi es amarrate le \VUgras{naves
a rotas} \textsuperscript{(31)} qui remonta le fluvio, e reguardante functionar
le \VUgras{gruas a mast} \textsuperscript{(32)}, qui leva, como plumas, enorme
cargas.

%%K t10-tit-5 sub
\VUsaut{4.0mm}
\VUpk{10.2pt}{40}{Le piloto.}
\VUsaut{3.0mm}

%%K t10-14-1 p
14. --- Io admirava le \VUgras{transatlanticos} \textsuperscript{(34)} exiente
del rada, con lor \VUgras{bandieras} \textsuperscript{(35)} in plen aere,
conducte per un habile \VUgras{piloto} \textsuperscript{(36)} qui meraviliosemente
sape sequer le \VUgras{canaletto} marcate per \VUgras{boias}
\textsuperscript{(40)} e \VUgras{balisas}, evitante le \VUgras{gabarras}
\textsuperscript{(41)} qui penosemente se dirige per lor unic vela quadrate. Io
distingueva sur le \VUgras{parte anterior} \textsuperscript{(37)} --- le prora
--- le \VUgras{cabinas} \textsuperscript{(39)} del secunde classe, e sur le
\VUgras{parte posterior} \textsuperscript{(38)} --- le puppa --- le salones pro
le passageros del prime classe.

%%K t10-15-1 p
15. --- Alcun vices io anque assisteva al cargamento de un \VUgras{chalana}
\textsuperscript{(42)} in le qual justo esseva amassate le mercantias del
\VUgras{magazin} \textsuperscript{(43)} e del \VUgras{station maritime}
\textsuperscript{(44)}, apportate per le numerose trainos del \VUgras{linea del
quai} \textsuperscript{(45)}.

%%K t10-tit-6 sub
\VUsaut{3.0mm}
\VUpk{10.2pt}{40}{Le placer del remos.}
\VUsaut{3.0mm}

%%K t10-16-1 p
16. --- Sovente io canotava in le \VUgras{docko a flottation}
\textsuperscript{(53)}, in le qual veni refugiar se le \VUgras{barcas}
\textsuperscript{(46)}, e il esseva un gaudio pro me manovrar le \VUgras{remo}
\textsuperscript{(47)}. Io ascendeva sur le \VUgras{ponte} \textsuperscript{(48)}
de un \VUgras{barca a vela} \textsuperscript{(49)}, io ascendeva, per
\VUgras{cordages} \textsuperscript{(52)}, sur le \VUgras{mast}
\textsuperscript{(51)} usque le \VUgras{antennas} \textsuperscript{(50)}, postea
io repeteva mi promenada \VUgras{in un yole} \textsuperscript{(54)} inter
pesante \VUgras{barcas de pisca} \textsuperscript{(55)}, ubi le \VUgras{retes}
\textsuperscript{(56)} se sicca.

%%K t10-17-1 p
17. --- Sur le \VUgras{quai} \textsuperscript{(58)} defilava ante me le plus
diverse \VUgras{mercantias} \textsuperscript{(59)}, contenite in \VUgras{cassas}
\textsuperscript{(60)} o in \VUgras{grande pacchettos} \textsuperscript{(61)}.
Illos formava le \VUgras{carga} \textsuperscript{(63)} del \VUgras{barcassas},
e potente \VUgras{gruas} \textsuperscript{(62)} los levava e los deponeva sur
\VUgras{camiones} \textsuperscript{(64)} durante que le \VUgras{carrettatores}
los declarava e pagava le taxas in le \VUgras{dogana} \textsuperscript{(65)}.

%%K t10-18-1 p
18. --- Finalmente, quando io arrivava al \VUgras{ponte giratori}
\textsuperscript{(66)} o al \VUgras{esclusa} \textsuperscript{(69)}, passante
ante le \VUgras{machina a dragar} \textsuperscript{(68)} que netta le
\VUgras{canal} \textsuperscript{(67)}, io disbarcava sur le quai presso le
\VUgras{semaphoro} \textsuperscript{(70)}.

%%K t10-19-1 p
19. --- Al altere latere de iste quai ja veni abordar le \VUgras{chalupas}
\textsuperscript{(71)} qui conduce a terra le officieros del esquadra. Il es un
admirabile spectaculo vider le marineros levar in cadentia lor remos, durante
que le barca portate per le \VUgras{unda} \textsuperscript{(72)}, facilemente
aborda le scala del disbarcatorio. Sur iste loco, on pote melio observar le
\VUgras{marea} e le movimento del \VUgras{fluxo} e \VUgras{refluxo} sur le
\VUgras{plagia} \textsuperscript{(73)}.

%%K t10-tit-7 sub
\VUsaut{3.0mm}
\VUpk{10.2pt}{40}{Le club.}
\VUsaut{3.0mm}

%%K t10-20-1 p
20. --- Pro retornar a casa, io usava le \VUgras{tramvia electric}
\textsuperscript{(74)} del qual le \VUgras{halto} \textsuperscript{(75)} es
presso le \VUgras{posto} ponite ante le club del officieros.

%%K t10-20-2 p
De sur le \VUgras{balcon} \textsuperscript{(76)} del club, ornate de
\VUgras{ancoras} \textsuperscript{(77)}, cannones e ballas, sovente io videva un
splendide reunion de officieros de marina : \VUgras{locotenentes}
\textsuperscript{(78)}, con lor spada, \VUgras{capitanos de artilleria}
\textsuperscript{(79)} e de \VUgras{infanteria} \textsuperscript{(80)} con lor
\VUgras{sabla}. Detra illes il habeva un \VUgras{marinero}
\textsuperscript{(81)} qui les apportava un \VUgras{telescopio}, quando illes
voleva distinguer lo que occurre lontano.

%%K t10-21-1 p
21. --- Io anque videva juvene \VUgras{sub-locotenentes} \textsuperscript{(82)}
recipiente le instructiones de lor \VUgras{commandante} \textsuperscript{(83)} o
del \VUgras{admiral} \textsuperscript{(84)}, durante que un \VUgras{quartiero-mastro}
\textsuperscript{(85)} attendeva pro portar los al nave.

%%K t10-22-1 p
22. --- De iste balcon, le vista es splendide : on domina tote le plagia con su
\VUgras{tentas} \textsuperscript{(86)} e \VUgras{cabinas} \textsuperscript{(87)},
e on vide, al hora del banio, le \VUgras{baniatores} \textsuperscript{(88)} qui
gaimente petula ante le \VUgras{casino} \textsuperscript{(89)}.

%%K t10-23-1 p
23. --- Al extremitate del plagia, le litore forma un parve \VUgras{peninsula}
\textsuperscript{(91)} \VUgras{rocose}, ubi le \VUgras{grande unda}
\textsuperscript{(92)} veni rumper se e sur le qual es construite un \VUgras{pharo}
\textsuperscript{(93)}, indicante de nocte le via al navigantes. Iste peninsula
communica con le terra solmente per un \VUgras{isthmo} multo stricte
\textsuperscript{(90)}.

%%K t10-tit-8 sub
\VUsaut{3.0mm}
\VUpk{10.2pt}{40}{Vista general del litores.}
\VUsaut{3.0mm}

%%K t10-24-1 p
24. --- Le \VUgras{scarpamento} \textsuperscript{(94)} se extende, findite per le
mar, que capriciosemente designa in illo un serie de \VUgras{bahias}
\textsuperscript{(95)} e \VUgras{promontorios} (capos), rendente lo multo
pittoresc.

%%K t10-24-2 p
Sur le \VUgras{plateau} \textsuperscript{(96)}, que domina le
\VUgras{stricto (de mar)} \textsuperscript{(98)}, il ha un \VUgras{forte}
\textsuperscript{(97)} multo importante e alcun « villas ».

%%K t10-25-1 p
25. --- Iste stricto separa del continente un \VUgras{insula}
\textsuperscript{(99)} multo periculose, circumferite de \VUgras{scolios}
\textsuperscript{(100)} e \VUgras{rumpe-undas}, ubi barcas e mesmo grande naves
naufraga alcun vices. Malgrado le promptitude del succursos, e le rapide arrivata
del \VUgras{barcas de salvamento}, iste \VUgras{naufragios} es terribile, e,
durante le \VUgras{tempestas} equinoctial, plure naves pereva con homines e
mercantias.

%%K t10-25-2 p
Malgrado iste vicinitate, \VUgras{le porto} es multo frequentate per le
marineros.
\VUsaut{6.0mm}
\VUfilet{22mm}
